% Source: https://artofproblemsolving.com/community/c3h2222788_fanzalist
% Stable identity: aops-topic-2222788-q01
\textbf{Tangent lengths.} Tangent segments from a point to a circle are equal. With $s=\dfrac{73+123+120}{2}=158$,
\[AE=AF=s-BC=35,\qquad BF=BD=s-CA=38,\qquad CD=CE=s-AB=85.\]

\textbf{Menelaus' theorem.} The line through $D$, $F$ and $G$ crosses the lines $BC$, $AB$, $CA$ of triangle $ABC$, so
\[\dfrac{AF}{FB}\cdot\dfrac{BD}{DC}\cdot\dfrac{CG}{GA}=1\qquad\Longrightarrow\qquad\dfrac{35}{38}\cdot\dfrac{38}{85}\cdot\dfrac{CG}{GA}=1\qquad\Longrightarrow\qquad\dfrac{CG}{GA}=\dfrac{85}{35}=\dfrac{17}{7}.\]
The line $DF$ meets the sides $BC$ and $AB$ at interior points, so it meets line $CA$ outside the segment $CA$; since $CG>GA$, the point $G$ lies beyond $A$.

\begin{asy}
usepackage("amsmath");
size(8.5cm);
defaultpen(fontsize(10pt));
pair B=(0,0), C=(123,0);
real ax=(73^2-120^2+123^2)/(2*123);
pair A=(ax,sqrt(73^2-ax^2));
pair D=(38,0);
pair F=B+(A-B)*38/73;
pair pE=C+(A-C)*85/120;
pair G=A+(A-C)*84/120;
pair I=(123*A+120*B+73*C)/(123+120+73);
real rr=abs(I-D);
draw(A--B--C--cycle,linewidth(0.8));
draw(circle(I,rr));
draw(A--G,dashed);
draw(D--G,blue);
dot(A);dot(B);dot(C);dot(D);dot(F);dot(pE);dot(G,red);dot(I);
label("$A$",A,NE);label("$B$",B,SW);label("$C$",C,SE);label("$D$",D,S);label("$F$",F,W);label("$E$",pE,NE);label("$G$",G,N,red);label("$I$",I,S);
label("$35$",(A+F)/2,W,fontsize(8pt));label("$38$",(F+B)/2,W,fontsize(8pt));
label("$38$",(B+D)/2,S,fontsize(8pt));label("$85$",(D+C)/2,S,fontsize(8pt));
label("$35$",(A+pE)/2,NE,fontsize(8pt));label("$85$",(pE+C)/2,NE,fontsize(8pt));
label("$84$",(G+A)/2,NE,red+fontsize(8pt));
\end{asy}

Let $GA=x$. Then $CG=x+120$ and $\dfrac{x+120}{x}=\dfrac{17}{7}$, so $7x+840=17x$ and $x=84$. Finally, $E$ lies on segment $AC$ with $AE=35$, so
\[EG=GA+AE=84+35=\boxed{119}.\]
